(1) 两圆轨道相切的转移椭圆必须以 \(r_1\) 为近地点、\(r_2\) 为远地点，故 \(a_t=(r_1+r_2)/2\)。轨道单位质量能量 \(\mathcal E=-\mu/(2a_t)\)，vis-viva公式 \(v^2=\mu(2/r-1/a_t)\) 给 \(\boxed{v_p=\sqrt{\mu(2/r_1-1/a_t)},\ v_a=\sqrt{\mu(2/r_2-1/a_t)}}\)。初末圆轨道速度分别 \(v_{c1}=\sqrt{\mu/r_1}\)、\(v_{c2}=\sqrt{\mu/r_2}\)。这里的椭圆两端是同一条无推力轨道的速度，不能分别按圆轨道公式计算。

(2) 因 \(a_t>r_1\)，第一次须加速 \(\Delta v_1=v_p-v_{c1}>0\)。因 \(a_t<r_2\)，椭圆远地点速度小于目标圆轨道速度，故第二次 \(\Delta v_2=v_{c2}-v_a>0\)。椭圆比角动量 \(h=\sqrt{\mu a_t(1-e^2)}=\sqrt{2\mu r_1r_2/(r_1+r_2)}\)，直接检验 \(r_1v_p=r_2v_a=h\)。每次脉冲引起单位质量机械能变化 \(\Delta\mathcal E=v_-\Delta v+(\Delta v)^2/2\)，两次相加恰等于 \(-\mu/(2r_2)+\mu/(2r_1)\)，但 \(\Delta v_1+\Delta v_2\) 不是该能量差的平方根。

(3) Kepler第三定律给椭圆周期 \(T_t=2\pi\sqrt{a_t^3/\mu}\)，从近地点到远地点历时半周期，故 \(\boxed{t_t=\pi\sqrt{a_t^3/\mu}}\)。初始圆周期 \(T_1=2\pi\sqrt{r_1^3/\mu}\)，所以 \(t_t/T_1=\frac12[(r_1+r_2)/(2r_1)]^{3/2}\)。若 \(r_2\to r_1\)，两脉冲趋零且转移时间趋半个圆周周期；若 \(r_2/r_1\to\infty\)，第一次速度增量趋 \((\sqrt2-1)v_{c1}\)，第二次趋零，但飞行时间按 \(r_2^{3/2}\) 发散。

(4) 理想火箭方程对每次脉冲给 \(m_{\rm before}/m_{\rm after}=\exp(\Delta v/v_e)\)。两次连续质量比相乘，得到 \(\boxed{m_0/m_f=\exp[(\Delta v_1+\Delta v_2)/v_e]}\)，推进剂质量 \(m_0-m_f\)。喷气反冲依赖速度增量而非单独的轨道势能差，且飞行器质量在喷气时改变，故不能把推进剂预算等同于轨道机械能增加量。若存在有限燃烧时间、重力损失、不同喷气速度或额外级间抛弃，必须重算。
